\documentclass[10pt]{beamer}

% Usando o tema Metropolis
\usetheme[block=fill]{metropolis} % ADICIONE block=fill
\usefonttheme{professionalfonts} % Para usar fontes personalizadas

% Carregando o pacote de fontes
\usepackage{beamerfontthemeslidestech}

% Configurações de formato de título
\pgfkeys{
  /slidestech/font/.cd,
  titleformat title=regular,      % Opções: regular, smallcaps, allsmallcaps, allcaps
  titleformat subtitle=regular, % Opções: regular, smallcaps, allsmallcaps, allcaps
  titleformat section=allcaps,    % Opções: regular, smallcaps, allsmallcaps, allcaps
  titleformat frame=smallcaps,    % Opções: regular, smallcaps, allsmallcaps, allcaps
}

% Metadados
\title{Exemplo do Tema de Fontes SlidesTech}
\subtitle{Demonstração Completa das Funcionalidades}
\author{Seu Nome}
\date{\today}

\begin{document}

% Slide de título
\begin{frame}
  \titlepage
\end{frame}

% Sumário
\begin{frame}{Sumário}
  \setbeamertemplate{section in toc}[sections numbered]
  \tableofcontents[hideallsubsections]
\end{frame}

% Seção 1
\section{Introdução}
\begin{frame}{Título do Frame em Small Caps}
  \begin{block}{Bloco de Exemplo}
    Este é um bloco com fonte em negrito no título, conforme configurado no código.
  \end{block}
  
  \begin{itemize}
    \item O título do frame está em \texttt{smallcaps}
    \item O título da seção está em \texttt{allcaps}
    \item O bloco tem título em negrito
  \end{itemize}
\end{frame}

\begin{frame}{Fontes Mono Espaçadas}
  \begin{columns}[T]
    \column{0.48\textwidth}
    Código exemplo:
    \texttt{function hello() \{ \\
    \quad return "World"; \\
    \}}
    
    \column{0.48\textwidth}
    Números em tabela:
    \begin{tabular}{r}
      \hline
      123456789 \\
      987654321 \\
      \hline
    \end{tabular}
  \end{columns}
\end{frame}

% Seção 2
\section{Personalização de Fontes}
\begin{frame}{Características das Fontes}
  \begin{itemize}
    \item \textbf{Roboto Light}: Fonte principal sem serifa
    \item \textit{Roboto Light Italic}: Versão itálica
    \item \textbf{Roboto}: Versão em negrito
    \item \textbf{\textit{Roboto Italic}}: Negrito itálico
    \item \texttt{Roboto Mono}: Fonte mono espaçada
  \end{itemize}
\end{frame}

\begin{frame}{Hierarquia de Títulos}
  \begin{block}{Título do Bloco}
    Este bloco demonstra a hierarquia de fontes.
  \end{block}
  
  \begin{alertblock}{Bloco de Alerta}
    Diferentes tipos de blocos mantêm a mesma formatação.
  \end{alertblock}
  
  \begin{exampleblock}{Bloco de Exemplo}
    A fonte do título do bloco é \texttt{\textbackslash normalsize} e \texttt{\textbackslash bfseries}.
  \end{exampleblock}
\end{frame}

% Seção 3
\section{Elementos Especiais}
\begin{frame}{Metadados e Legendas}
  \begin{figure}
    \rule{0.8\textwidth}{0.4\textwidth}
    \caption{Esta é uma legenda com fonte small e nome em negrito}
  \end{figure}
  
  \begin{itemize}
    \item Autor e data aparecem com fonte \texttt{small}
    \item Números de página com fonte \texttt{scriptsize}
  \end{itemize}
\end{frame}

\begin{frame}{Referências Bibliográficas}
  \begin{thebibliography}{99}
    \bibitem{exemplo} Autor Exemplo.
    \newblock {\bfseries Título do Artigo}.
    \newblock {\normalfont Journal Name}, 2024.
    \newblock {\small Nota: Demonstração de formatação}.
  \end{thebibliography}
\end{frame}

\begin{frame}[standout]
  \centering
  \Large\textbf{Slide de Destaque (Standout)}
  
  \normalsize Com fonte grande e negrito
\end{frame}

% Seção 4
\section{Demonstração de Formatos}
\begin{frame}{Comparação de Formatos de Título}
  \framesubtitle{Título do frame em smallcaps, subtítulo em smallcaps}
  
  Aqui podemos ver:
  \begin{enumerate}
    \item Título da seção: \MakeUppercase{formato allcaps}
    \item Título do frame: \textsc{formato smallcaps}
    \item Subtítulo do frame: \textsc{também smallcaps}
  \end{enumerate}
\end{frame}

\begin{frame}{Tabela com Números Mono Espaçados}
  \begin{table}
    \caption{Exemplo de tabela com números}
    \begin{tabular}{lr}
      \hline
      Item & Quantidade \\
      \hline
      Item A & 12345 \\
      Item B & 67890 \\
      Item C & 11111 \\
      \hline
    \end{tabular}
  \end{table}
  
  Os números na tabela usam fonte mono espaçada automaticamente.
\end{frame}

\end{document}
